% BHU PYQ -- Manually transcribed & error-corrected
% Paper: BCHB-327 : Law & Practice of Insurance
% Exam: B.Com. (Hons.) (Semester VI) Examination, 2015-16

\documentclass[11pt,a4paper]{article}
\usepackage[a4paper,top=2.5cm,bottom=2.5cm,left=2.8cm,right=2.8cm,headheight=14pt]{geometry}
\usepackage{amsmath,amssymb}
\usepackage[T1]{fontenc}
\usepackage[utf8]{inputenc}
\usepackage{lmodern,microtype,fancyhdr,enumitem,hyperref,lastpage,parskip}

\hypersetup{
    colorlinks=true,
    urlcolor=black,
    linkcolor=black,
    pdftitle={BCHB-327: Law \& Practice of Insurance -- BHU B.Com. (Hons.) Sem VI 2015-16}
}

\pagestyle{fancy}
\fancyhf{}
\renewcommand{\headrulewidth}{0.4pt}
\renewcommand{\footrulewidth}{0.4pt}
\fancyhead[L]{\small   \textit{Banaras Hindu University}}
\fancyhead[R]{\small   \textit{BCHB-327: Law \& Practice of Insurance}}
\fancyfoot[C]{\small Page  \thepage\ of \pageref{LastPage}}


\newlist{parts}{enumerate}{2}
\setlist[parts,1]{label=(\alph*),leftmargin=*,itemsep=5pt,topsep=3pt}
\setlist[parts,2]{label=(\roman*),leftmargin=*,itemsep=3pt,topsep=2pt}

\newcommand{\pts}[1]{\hfill   \textit{[#1]}}


\begin{document}


\begin{center}
  {\large\bfseries BANARAS HINDU UNIVERSITY}\\[2pt]
  {
\normalsize B.Com. (Hons.) --- Semester VI Examination, 2015-16}\\[6pt]
  \rule{0.8\linewidth}{0.5pt}\\[6pt]
  {\large\bfseries Elective Group --- Banking \& Insurance}\\[4pt]
  {\large\bfseries Paper No. BCHB-327: Law \& Practice of Insurance}\\[4pt]
  {
\normalsize Time: 3 Hours\hfill Maximum Marks: 70}
\end{center}

\rule{\linewidth}{0.4pt}

\smallskip



\noindent   \textit{Instructions: Attempt all questions. All questions carry equal marks (14 marks each).}

\bigskip




\noindent   \textbf{Question 1.} \\
What do you understand by Insurance? Describe the provisions relating to registration under the Insurance Act, 1938.\pts{14}

\centerline{   \textbf{OR}}

Throw light on the main provisions of the Insurance Act, 1938.\pts{14}

\medskip\rule{\linewidth}{0.2pt}




\noindent   \textbf{Question 2.} \\
What do you mean by GIC? What are its main functions? Explain.\pts{14}

\centerline{   \textbf{OR}}

Discuss the main provisions of the LIC Act and also throw light on the organizational structure of the Life Insurance Corporation of India.\pts{14}

\medskip\rule{\linewidth}{0.2pt}

\pagebreak




\noindent   \textbf{Question 3.} \\
Who is the 'Authority' under IRDA? Express the various powers and functions of the Authority.\pts{14}

\centerline{   \textbf{OR}}

What do you understand by an insurance agent? Critically examine the provisions related to them.\pts{14}

\medskip\rule{\linewidth}{0.2pt}




\noindent   \textbf{Question 4.} \\
``Interests of Policy Holders'' is one of the important functions of Insurance. Comment with reference to IRDA.\pts{14}

\centerline{   \textbf{OR}}

Discuss the provisions relating to Micro Insurance under the IRDA Act.\pts{14}

\medskip\rule{\linewidth}{0.2pt}




\noindent   \textbf{Question 5.} \\
Discuss the steps involved in claims settlement.\pts{14}

\centerline{   \textbf{OR}}

Write notes on any two of the following:

\begin{parts}
  \item History of Insurance\pts{7}
  \item Underwriting in Insurance\pts{7}
  \item Re-insurance\pts{7}
\end{parts}

\vfill
\rule{\linewidth}{0.4pt}

\begin{center}\small
\href{https://bhu-exam.vercel.app/}{Click here to visit our website}\\[4pt]
\href{https://me-aryan.vercel.app/}{Click here to visit our contributor}\\[4pt]
\href{https://quashan-me.vercel.app/}{Click here to visit our contributor}
\end{center}

\end{document}
